\section{Conclusion}
As the adoption of Retrieval-Augmented Generation (RAG) accelerates in privacy-critical domains, standardizing the balance between semantic utility and data security is paramount. In this work, we introduced SHADE, a quantitative metric for measuring semantic reconstruction risk in distributed knowledge bases. We subsequently proposed CoShard, a graph-based partitioning algorithm explicitly designed to optimize the SHADE-coherence trade-off. 

Through rigorous theoretical analysis and empirical adversarial simulation, we established a No-Free-Lunch theorem for semantic sharding: maximizing utility inherently maximizes disclosure risk. However, by explicitly penalizing intra-shard semantic density, CoShard successfully traces the Pareto frontier, demonstrating absolute dominance over state-of-the-art privacy baselines (VertiSplitRAG, DPVoteRAG) and utility baselines (GraphRAG, K-Means). CoShard enables system administrators to enforce strict, mathematically sound privacy guarantees without suffering catastrophic utility degradation.
